\subsection{Kanonische Einheitsvektoren}
\label{subsec:vector_algebra_canonical_unit_vectors}

Einheitsvektoren können in jede beliebige Richtung zeigen. Besonders
naheliegend sind Einheitsvektoren, die entlang der Koordinatenachsen zeigen.

\begin{DefinitionBox}{Kanonische Einheitsvektoren}{vector_algebra_canonical_unit_vectors}
  Die \emph{kanonischen Einheitsvektoren}
  \EnglishTerm{canonical unit vectors} sind die Einheitsvektoren (Länge 1) entlang der
  positiven Koordinatenachsen:
  \begin{itemize}[leftmargin=1.5em,itemsep=0.9em]
    \item Im $\R^2$:
      \par\smallskip
      \begin{minipage}[c]{0.47\linewidth}
        \[
          \begin{gathered}
            \vec e_1=\begin{pmatrix}1\\0\end{pmatrix},
            \qquad
            \vec e_2=\begin{pmatrix}0\\1\end{pmatrix},\\[0.5em]
            \text{paarweise orthogonal}:
            \vec e_1\perp\vec e_2.
          \end{gathered}
        \]
      \end{minipage}\hfill
      \begin{minipage}[c]{0.47\linewidth}
        \centering
        \CachedTikzFigure{content/resources/chapters/vector_algebra/tikz/canonical_unit_vectors_2d.tex}
      \end{minipage}

    \item Im $\R^3$:
      \par\smallskip
      \begin{minipage}[c]{0.47\linewidth}
        \[
          \begin{gathered}
            \vec e_1=\begin{pmatrix}1\\0\\0\end{pmatrix},
            \quad
            \vec e_2=\begin{pmatrix}0\\1\\0\end{pmatrix},
            \quad
            \vec e_3=\begin{pmatrix}0\\0\\1\end{pmatrix},\\[0.5em]
            \text{paarweise orthogonal:}\\[0.25em]
            \vec e_1\perp\vec e_2,
            \quad
            \vec e_1\perp\vec e_3,
            \quad
            \vec e_2\perp\vec e_3.
          \end{gathered}
        \]
      \end{minipage}\hfill
      \begin{minipage}[c]{0.47\linewidth}
        \centering
        \CachedTikzFigure{content/resources/chapters/vector_algebra/tikz/canonical_unit_vectors_3d.tex}
      \end{minipage}
  \end{itemize}
\end{DefinitionBox}

Das Fremdwort \emph{kanonisch} wird in vielen Fachgebieten für eine besonders
naheliegende oder ausgezeichnete Standardwahl verwendet.

Vektoren werden daher nicht nur als Koordinatenspalten, sondern auch mithilfe
kanonischer Einheitsvektoren dargestellt. Diese insbesondere in der
Technischen Mechanik wichtige Darstellungsform macht die Anteile in den
einzelnen Koordinatenrichtungen unmittelbar sichtbar.

\begin{ExampleBox}{Darstellung mit kanonischen Einheitsvektoren}{vector_algebra_canonical_unit_vector_decomposition}
  \begin{enumerate}[label=\alph*),leftmargin=1.7em,itemsep=1.3em]
    \item
        $\displaystyle
        \begin{pmatrix}2\\3\end{pmatrix}
        =\begin{ClozeContinuation}[28]{2}
          \begin{pmatrix}2\\0\end{pmatrix} + \begin{pmatrix}0\\3\end{pmatrix}
          = 2 \begin{pmatrix}1\\0\end{pmatrix} +3 \begin{pmatrix}0\\1\end{pmatrix}
          = 2\vec e_1+3\vec e_2
        \end{ClozeContinuation}.$

    \item
      $\displaystyle
        \begin{pmatrix}1\\2\\3\end{pmatrix}
        =\begin{ClozeContinuation}[28]{3}
          \begin{pmatrix}1\\0\\0\end{pmatrix}
          +\begin{pmatrix}0\\2\\0\end{pmatrix}
          +\begin{pmatrix}0\\0\\3\end{pmatrix}
          =1\begin{pmatrix}1\\0\\0\end{pmatrix}
          +2\begin{pmatrix}0\\1\\0\end{pmatrix}
          +3\begin{pmatrix}0\\0\\1\end{pmatrix}
          =1\vec e_1+2\vec e_2+3\vec e_3
        \end{ClozeContinuation}.$
  \end{enumerate}
  \textbf{Erkenntnis:} Die Gewichtungsfaktoren der kanonischen Einheitsvektoren $\vec{e}_i$
  entsprechen den Komponenten des Vektors.
\end{ExampleBox}

\begin{KeyPointBox}{Zerlegung in kanonische Einheitsvektoren}{vector_algebra_canonical_unit_vector_representation}
  Jeder Vektor ist die gewichtete Summe der kanonischen Einheitsvektoren:
  \[
    \begin{array}{@{}r@{:\qquad}c@{\quad=\quad}l@{}}
      \R^2
        & \begin{pmatrix}a_1\\a_2\end{pmatrix}
          \vphantom{\begin{pmatrix}a_1\\a_2\\a_3\end{pmatrix}}
        & \Cloze{10cm}{a_1\vec e_1+a_2\vec e_2},\\[1.4em]
      \R^3
        & \begin{pmatrix}a_1\\a_2\\a_3\end{pmatrix}
        & \Cloze{10cm}{a_1\vec e_1+a_2\vec e_2+a_3\vec e_3},\\[1.4em]
      \R^n
        & \begin{pmatrix}a_1\\a_2\\\vdots\\a_n\end{pmatrix}
        & \Cloze{10cm}{a_1\vec e_1+a_2\vec e_2+\dots+a_n\vec e_n}.
    \end{array}
  \]
\end{KeyPointBox}
